\documentclass[12pt,english,xcolor=dvipsnames,aspectratio=1610]{beamer}

\usepackage[utf8]{inputenc}
\usepackage{color}
\usepackage[sort,round]{natbib}
\usepackage{booktabs}
\usepackage{graphicx}
\usepackage{threeparttable}
\usepackage{cancel}
\usepackage{lscape}
\usepackage{tikz}
\usetikzlibrary{calc}
\beamertemplatenavigationsymbolsempty

\setbeamercolor{title}{fg=Blue,bg=}
\setbeamercolor{frametitle}{fg=Blue,bg=}
\setbeamercolor{framesubtitle}{fg=Blue,bg=}
\setbeamercolor{caption name}{fg=Blue, bg=}

\newcommand{\myitem}{\item[\textbullet]}
\usetheme{boxes}
\usecolortheme{default}

\title{Beneath the Surface: Investigating Anti-Competitive Behavior in Brazilian Public Procurement}
\author{Sergio Galletta, Darcio Genicolo-Martins, Tommaso Giommoni}
\institute{ETH Zürich, INSPER, University of Amsterdam}
\date{16 October 2024}

\begin{document}

\begin{frame}
    \titlepage
\end{frame}

% Slide 1: Motivation
\begin{frame}{Motivation}
    \begin{itemize}
        \item Integrity and efficiency in procurement are vital for the effective use of public funds, economic development, and public trust (Bosio et al, 2022).\vspace{2mm}
        \item Public procurement is a significant component of the global GDP (12\% overall and up to 20\% in developing countries).\vspace{2mm}
        \item Anti-competitive behavior (collusion, favoritism) increases costs, reduces service quality, limits innovation.\vspace{2mm}
        \item Detecting anti-competitive behavior in procurement is difficult due to its hidden nature and the complexity of the processes involved. \vspace{2mm}
    \end{itemize}
\end{frame}

\begin{frame}{Introduction - The Paper}
\begin{itemize}
\item \textbf{Research topic:} Detect signals of anti-competitive behavior in public auctions in São Paulo, Brazil. \\ \vspace{2mm} \pause
\item \textbf{Data:} 3.8 million procurement auctions from 2010 to 2019, including bids and firm-level information. \\ \vspace{2mm} \pause
\item \textbf{Methods:} \\ \vspace{2mm}
\begin{itemize}
\item [-] We use a \textbf{Regression Discontinuity Design (RDD)}, exploiting small differences between narrowly winning and narrowly losing bids (Kawai et al., 2023 - RESTUD). \\ \vspace{2mm}
\item [-] Winning is nearly random; if \textbf{incumbency} or \textbf{backlog} spikes at the threshold, it signals potential anti-competitive behavior. \pause \\ \vspace{2mm}
\end{itemize}
\item \textbf{Results:} Some evidence of potential systemic collusion, yet more to be done! 
\end{itemize}
\end{frame}

\begin{frame}{Anti-competitive Behavior}
    \begin{itemize}
        \item Different reasons for anti-competitive behavior when the public sector is involved\vspace{2mm}
        \item \textbf{Collusion:} Agreements among firms can lead them to overcharge for goods and services.\vspace{2mm}
        \item \textbf{Corruption:} Government officials may award contracts to politically connected firms over more efficient ones\vspace{2mm}
    \end{itemize}
\end{frame}

\begin{frame}
\frametitle{Identifying Anti-Competitive Behaviors in Public Procurements}

\begin{itemize}
    \item Commonly identified strategies \vspace{2mm}
    \begin{enumerate}
    \item \textbf{Incumbency priority}: Winners are more likely to be incumbents\vspace{2mm}
    \item \textbf{Bid rotation}: Winners are likely to have lower backlog
    \end{enumerate}
     \vspace{2mm}\pause
    \item The \textbf{problem} is that there are \textbf{non-collusive cost-based explanations} for these allocation patterns\vspace{2mm}
    \begin{itemize}
        \item Bid rotation can arise under competition if \textbf{marginal costs increase with backlog}.\vspace{2mm}
        \item Incumbency advantage can be explained by \textbf{cost asymmetries} among competitive firms or by \textbf{learning-by-doing}.
    \end{itemize}
    \end{itemize}
\end{frame}

\begin{frame}
\frametitle{Identifying Anti-Competitive Behaviors in Public Procurements}

\begin{itemize}
    \item Kawai et al.'s (2023) approach hinges on the idea that \textbf{closely matched bids imply similar cost structures and winning probabilities.}\vspace{2mm}
    \item Discrepancies in firms' previous outcomes could indicate anti-competitive behavior.\vspace{2mm}\pause
    \item By using a \textbf{Regression Discontinuity Design (RDD)}, one can exploit small differences between narrowly winning and narrowly losing bids.\vspace{2mm}
    \item \textbf{Key Assumption:} Firms that bid just above or just below the winning margin \textbf{ should be very similar in terms of cost structures}.\vspace{2mm}\pause
    \item We are interested in \textbf{1) Incumbency Status} and \textbf{2) Backlog} of firms bidding around the cutoff. 
    \end{itemize}
\end{frame}


\begin{frame}{Regression Discontinuity Design (RDD)}
    \[
    \beta = \lim_{\epsilon \to 0^+} E[x_{i,t} | \Delta_{i,t} = \epsilon] - \lim_{\epsilon \to 0^-} E[x_{i,t} | \Delta_{i,t} = \epsilon]
    \]
    where:
    \begin{itemize}
        \item \( \Delta_{i,t} \): Difference between the bid of firm \( i \) and the most competitive alternative bid at time \( t \).
        \item \( x_{i,t} \): Outcome of interest (incumbency or backlog).
    \end{itemize}
\end{frame}

\begin{frame}{Institutional Background}
    \begin{itemize}
            \item São Paulo is the wealthiest and most populous state (46 million inhabitants) in Brazil, contributing over 30\% of the country's GDP.\vspace{2mm}
            \item Major industries include manufacturing, finance, agriculture, and services. \vspace{2mm}\vspace{2mm}\pause
        
            \item São Paulo has one of the largest public procurement markets in Brazil, handling billions in public contracts annually. \vspace{2mm}
             \item The large scale of procurement creates opportunities for anti-competitive behavior (e.g., a detected school meal cartel was fined over R$340 million/US$60 million).\vspace{2mm}
    \end{itemize}
\end{frame}


\begin{frame}{Institutional Background}
    \begin{itemize}
        \item São Paulo uses the \textbf{Bolsa Eletrônica de Compras (BEC)} platform for public procurement.\vspace{2mm}
        \item Two types of tenders:\vspace{2mm}
        \begin{itemize}
            \item \textbf{\textcolor{red}{Sealed bid tenders (convite)}}: Suitable for smaller contracts (under R\$176,000 or - approximately US\$31,000).\vspace{2mm}
            \item \textbf{Multi-round descending auctions (pregão)}: For larger contracts, includes a reverse auction phase.\vspace{2mm}
        \end{itemize}
    \end{itemize}
\end{frame}


% Slide 3: Data
\begin{frame}{Data}
    \begin{itemize}
        \item Data from the BEC platform (2009–2019).\vspace{2mm}
        \item \textbf{3.8 million records}, covering 19,000 firms and 1,344 public buyer units.\vspace{2mm}
        \item Key variables:\vspace{2mm}
        \begin{itemize}
            \item \textbf{Incumbency}: Incumbency: Whether a firm won the previous auction it participated in\vspace{2mm}
            \item \textbf{Ratio Won}: Share of auctions won in the previous months \vspace{2mm}
            \item \textbf{Backlog}: Total amount of work a firm has accumulated in the previous months.\vspace{2mm}
            \item \textbf{Margin of Victory}: Percentage difference between the lowest and second-lowest bids.\vspace{2mm}
            \item \textbf{Firm characteristics}: age, legal status, and more to be collected
        \end{itemize}
    \end{itemize}
\end{frame}

\begin{frame}{Data}
    \begin{figure}
        \includegraphics[width=0.8\textwidth]{total_amount_awarded_per_year (1).pdf}
    \end{figure}
\end{frame}

\begin{frame}{Graphical Results - Incumbency advantage}
    \begin{figure}
        \includegraphics[width=0.8\textwidth]{incumbent.png}
    \end{figure}
\end{frame}


\begin{frame}{Graphical Results - Ratio Won}
    \begin{figure}
        \includegraphics[width=0.8\textwidth]{share_won.png}
    \end{figure}
\end{frame}

\begin{frame}{Graphical Results - Backlog 120 days}
    \begin{figure}
        \includegraphics[width=0.8\textwidth]{backlog.png}
    \end{figure}
\end{frame}


\begin{frame}{Numerical Estimates}
\begin{center}
\begin{tabular}{l*{6}{c}} \hline\hline                                                  
&Incumbent   & Ratio Won    & Backlog 120     \\  
&  (1)                  & (2)                   & (3)  \\ 
\hline \\ 
 Winning Auction           &     0.033***  & 0.016*** &  0.027*** \\
        & (0.004) & (0.004) & (0.009)   \\   \\ 
Observations    &   146,999  &  76,572 & 89,307 \\  
Bandwidth      &  0.01        & 0.01 & 0.01 \\
Mean DV & 0.672 & 0.494 & 0 \\
SD DV & 0.469 & 0.326 & 1 \\
 \hline \hline \end{tabular}      
\end{center}
\end{frame}

\begin{frame}{Balance checks}
\begin{center}
\begin{tabular}{l*{6}{c}} \hline\hline                                                  
&Firm Age   & Limited liability    & Buyer/vendor     \\  
&   &     & same municipality     \\  
&  (1)                  & (2)                   & (3)  \\ 
\hline \\ 
 Winning Auction           &     0.026  & 0.002 &  -0.003 \\
        & (0.077) & (0.004) & (0.002)   \\   \\ 
Observations    &   208,732  &  208,984& 208,984 \\  
Bandwidth      &  0.01        & 0.01 & 0.01 \\
Mean DV & 10.82 & 0.571 & 0.200 \\
SD DV & 10.43 & 0.494 & .400 \\
\hline \hline \end{tabular}      
\end{center}
\end{frame}


\begin{frame}{Future analysis}
 \begin{itemize}
        \item Balance checks on other firms' characteristics
        \item Analysis by kind of good, buyer, year, municipality\vspace{2mm}
        \item Estimate the $\beta$s and check the effect of exogenous shocks on anti-competitive behavior:\vspace{2mm}
        \begin{itemize}
            \item Effects of exogenous audits (i.e., random audits)\vspace{2mm}
            \item Organized crime effect (i.e., geographical expansion of PCC—Primeiro Comando da Capital)
        \end{itemize}
    \end{itemize}
\end{frame}



% Slide 7: Conclusion
\begin{frame}{Conclusion}
    \begin{itemize}
        \item Narrowly winning firms show higher incumbency and larger backlogs
        \item They are similar in age and legal structure 
        \item Overall, no evidence of bid rotation
        \item Results point to the potential presence of systemic anti-competitive behavior
    \end{itemize}
\end{frame}


\begin{frame}
\begin{center}
    \Huge Thank you!
\end{center}
\end{frame}

\end{document}
